\section*{Capítulo \ref{ch:g5:division} --- Divisão e múltiplos}

\begin{solution}{exo:g5:division:1}
$851 \div 23$: \omterm{def:g5:division:multiple}{múltiplos} de $23$: $23, 46, 69, 92, 115, 138, 161,
184, 207$. Então $85 \div 23$: $3$ vezes ($69$), resto $16$;
derrubar o $1$: $161 \div 23 = 7$ exatamente. quociente $37$,
resto $0$ ($23 \times 37 = 851$).

$704 \div 32$: $70 \div 32$: $2$ vezes ($64$), resto $6$; trazer
para baixo $4$: $64 \div 32 = 2$. quociente $22$, resto $0$.
\end{solution}

\begin{solution}{exo:g5:division:2}
$9 \div 2 = 4.5$; \quad $27 \div 4 = 6.75$; \quad
$33 \div 5 = 6.6$; \quad $21 \div 8 = 2.625$.
\end{solution}

\begin{solution}{exo:g5:division:3}
$54 \div 4 = 13.5$: cada um paga $13.50$.
\end{solution}

\begin{solution}{exo:g5:division:4}
$10 \div 3 = 3.333\dots $: o dígito $3$ repete para sempre --- o
a divisão nunca para. Arredondado para centésimo: $3.33$.
\end{solution}

\begin{solution}{exo:g5:division:5}
múltiplos de $7$ até $70$: $7, 14, 21, 28, 35, 42, 49, 56, 63,
70$. Divisores de $18$: $1, 2, 3, 6, 9, 18$. Divisores de $24$:
$1, 2, 3, 4, 6, 8, 12, 24$.
\end{solution}

\begin{solution}{exo:g5:division:6}
$470$: por $2$ (termina em $0$), por $5$, por $10$; dígito soma $11$: não por
$3$.

$735$: termina em $5$: por $5$, não por $2$ nem $10$; dígito soma $15$: por
$3$.

$8\,001$: ímpar; dígito soma $9$: por $3$ apenas.

$1\,314$: par; dígito soma $9$: por $2$ e por $3$ (portanto por $6$),
não por $5$.
\end{solution}

\begin{solution}{exo:g5:division:7}
Divisível por $3$ e $5$ significa divisível por $15$: entre $60$ e
$80$, os números $60$, $75$.
\end{solution}

\begin{solution}{exo:g5:division:8}
$7.2 \div 6 = 1.2$: cada peça mede $1.2$ m (verifique:
$1.2 \times 6 = 7.2$).
\end{solution}

\begin{solution}{exo:g5:division:9}
$1\,000 = 24 \times 41 + 16$: quarenta e um sacos cheios, $16$ bolinhas de gude
esquerda.
\end{solution}

\begin{solution}{exo:g5:division:10}
$7 \div 5 = 1.4$: cada um dos $5$ amigos recebem $1.4$ bares ---
possível cortando as barras em décimos. Para $5 \div 7$: a divisão
dá $0.714285\dots $, que nunca para --- nenhuma fração decimal exata
existe; apenas o fração $\frac57$ nomeia exatamente.
\end{solution}

\begin{solution}{exo:g5:division:11}
Dígito soma: $5 + 2 + d = 7 + d $. Divisível por $3$ quando $7 + d $ é
$9$, $12$ ou $15$: $ d = 2$, $5$ ou $8$. O menor é $ d = 2$
(dando $522 = 3 \times 174$).
\end{solution}
